\begin{abstract}
\noindent
The Yang--Mills workbench of KokunoYumeto, published as the Zenodo record 10.5281/zenodo.22883643, \emph{SU(2) lattice Yang--Mills: vacuum, energy, spectral bounds and continuum investigations}, studies the Kogut--Susskind Hamiltonian of $SU(2)$ lattice gauge theory on open cubic boxes. The record holds 43 files from five editions, dated 8 to 21 September 2026. This paper says which files are current, what each document claims and on which regime, and what is proved, with proofs and with the programs that check the constants.

The central result is an explicit lower bound on the spectral gap above the ground state, on the whole gauge-invariant space, uniform in the box size and the lattice spacing, for bare coupling $g^2\ge3.6836$ in the workbench's normalisation; the threshold rests on computed coefficient bounds of orders 3 to 5. With order-2 inputs recomputed here, the same argument gives a gap of at least $1.52\kappa$ for $g^2\ge3.9781$, which includes the workbench's benchmark $g^2=4$. Its key combinatorial input is sharp: on any graph, gauge invariance forces the spins of a Peter--Weyl block to satisfy $\sum_fj_f\ge g\,j_e$ with $g$ the girth, and the girth is the best uniform constant. From the same control of the vacuum the paper derives that the gap persists at complex coupling, so that the ground state is analytic in the coupling on a disk whose radius does not depend on the box; that the gap on the full, not gauge-invariant, space is at least a quarter of the physical one; and that for $g^2\ge13$ the gap lies in $[2.9047,3.00024]\kappa$. For the heat flow of plaquettes it re-derives every constant of the box-independent heat remainder, improves its prefactor from $67896/169$ to $2829/13$, and sharpens the bounds for relaxation into the complementary physical space.

The paper states the proved limits of the global perturbative route and of the workbench's standing continuum path, the typed maps by which the workbench brings the Jacobian-map, Navier--Stokes and $S^6$ material into the Hamiltonian setting, and the directions that remain open. All results hold at a fixed lattice spacing and strong bare coupling; the workbench does not claim a continuum mass gap, and nothing here claims one.
\end{abstract}

\section*{Introduction}\phantomsection\label{sec:about}
\addcontentsline{toc}{section}{Introduction}

\subsection*{The problem and this record}

The Clay Mathematics Institute's problem asks for a quantum Yang--Mills theory on $\R^4$ with a mass gap \cite{JaffeWitten}. Lattice gauge theory approaches it through regulated models, and at strong coupling mass gaps are classical in the Euclidean setting \cite{OsterwalderSeiler}. The workbench works in the Hamiltonian setting of Kogut and Susskind \cite{KogutSusskind}. It aims at explicit bounds that are uniform in the volume, hold on the full physical space and have closed coupling thresholds, together with what such bounds make available: heat correlations, relaxation into the complementary space, and analyticity in the coupling. Its material spans five Zenodo editions and several hundred pages.

The workbench is research carried out with AI systems, which its author used to formulate and test many hypotheses quickly. Some of its directions are finished theorems, some are programmes with proved parts, and some are proposals; not every direction is expected to hold. An unfinished programme is still mathematics, and its proved parts stand on their proofs whatever becomes of the rest. This paper states what holds, with proofs or the exact location of a proof, what fails and in which scope, and how the workbench connects with the author's other programmes. The audit behind it had limited compute and did not examine every direction; where the paper gives a view rather than a proof, the view is labelled as mine and its reasons are given.

\subsection*{Main results}

The Hamiltonian is $H_L=\kappa K+\kappa\xi(2M_L-S)$ on the open box $\{-L,\dots,L\}^3$, with $K$ the sum of the link Casimirs, $S$ the sum of the plaquette traces, $M_L$ the number of plaquettes, $\kappa=2g^2/a$ and $\xi=1/(4g^4)$ (Section~\ref{sec:problem}). The gap $\Delta_L$ is the distance from the ground energy to the rest of the spectrum on the gauge-invariant space; at $\xi=0$ it equals $3\kappa$.

\begin{introthm}[Theorem~\ref{thm:ym01}, Corollary~\ref{cor:ym-orderN}]
If the workbench's computed bounds on the local norms of the vacuum coefficients of orders 1 to 5 hold, then for every box $L\ge2$ and every spacing $a>0$, $\Delta_L\ge\kappa\,d_5(\xi)$ whenever $g^2\ge3.6836$, with an explicit function $d_5$ decreasing from $3$; in particular $\Delta_L\ge1.5453\kappa$ on that range and $\Delta_L>1.9068\kappa$ for $g^2\ge4$. With the order-1 and order-2 inputs verified here, $\Delta_L\ge1.52054\kappa$ for all $g^2\ge3.97807$.
\end{introthm}

The proof writes the ground state as $e^v/\|e^v\|$, solves the resulting equation for $v$ order by order, controls the remainder in local Fourier-algebra norms that do not grow with the box, and returns to the eigenvalue equation. Its combinatorial input is the following inequality, which the workbench proves for the cubic lattice and which holds on every graph.

\begin{introthm}[Lemma~\ref{lem:casimir}, Remark~\ref{rem:perlink}]
Let the links form a finite graph in which every cycle has length at least $g$. On every Peter--Weyl block $j$ that contains a nonzero gauge-invariant vector, $\sum_fj_f\ge g\,j_e$ for every link $e$. The girth is the best constant uniform over the links, and the exact constant for each link is determined by the shortest cycles and closed walks through its endpoints.
\end{introthm}

For the cubic box the inequality gives $c_j\ge3$ for the Casimir of every nonconstant physical block, which is what makes the spectral return work. The same control of the vacuum gives more than a lower bound.

\begin{introthm}[Propositions~\ref{prop:complexgap} and~\ref{prop:gapupper}, Corollary~\ref{cor:fullspace}]
Under the inputs of Theorem~I through order $N$:
(i) for complex coupling $|\zeta|\le\alpha_N$ the ground-state eigenvalue of $H_L(\zeta)$ is algebraically simple and every other eigenvalue has real part at least $\kappa\,d_N(|\zeta|)$ above it, so the ground-state projection is analytic in the coupling on a disk whose radius does not depend on the box;
(ii) the gap on the whole space $L^2$ is at least $\frac14\kappa\,d_N(\xi)$, which is sharp at $\xi=0$;
(iii) for every box, $\Delta_L$ is bounded above by an explicit plaquette quotient, and for $g^2\ge13$ the gap lies in $[2.9047,3.00024]\kappa$; at each fixed $L$, $\Delta_L$ has no first-order term in $\xi$ at $\xi=0$.
\end{introthm}

The second manuscript of the latest edition uses the vacuum for heat correlations of the centred plaquette states $r_p$.

\begin{introthm}[Theorem~\ref{thm:heat}]
For every $L\ge2$, $\tau\ge0$ and $|\xi|<1/55$, the heat correlations $\widehat C(\tau;\xi)$ of the plaquette states differ from their Taylor polynomial of degree $4$ in $\xi$ by at most $\frac{2829}{13}e^{-13\tau/8}(55|\xi|)^6/(1-(55|\xi|)^2)$ in row norm.
\end{introthm}

The workbench states this with the prefactor $67896/169$. The relaxation bounds of Section~\ref{ssec:complement} then show, under the inherited inputs of Proposition~\ref{prop:bench}, that at $g^2\ge13$ the plaquette family loses less than a fraction $0.000237$ of its energy when it is allowed to relax into the whole complementary physical space.

\subsection*{What the results mean}

At strong coupling the workbench obtains, for the $SU(2)$ Hamiltonian on open boxes, a spectral gap with explicit constants that is uniform in the volume and holds on the entire physical space. The proof rests on two structural facts: gauge invariance forces the Casimir of every nonconstant physical block to be at least $3$, because the girth of the cubic lattice is $4$; and the vacuum, written as an exponential, can be controlled in local norms that do not grow with the box. The same control gives box-independent heat correlations, near-complete relaxation of the plaquette family, analyticity of the ground state in the coupling on a box-independent disk, and an upper bound that pins the gap within $0.03\%$ of $3\kappa$ at $g^2\ge13$.

These are statements at fixed lattice spacing and strong bare coupling. The continuum problem needs $a\to0$ with $g\to0$, the opposite regime, and the workbench's standing continuum path leaves the domain of these estimates after finitely many steps (Proposition~\ref{prop:path}). What would connect the two regimes is a bound of this kind that persists along a path to weak coupling; Section~\ref{sec:directions} says which ingredients of the present proof would have to be extended. The workbench also pursues the opposite construction: sequences driven by the $S^6$ period data in which box, spacing and coupling vary together, aimed at models that are gapped at every regulator but whose continuum limit is gapless.

\subsection*{The bridges}

Three branches of the workbench bring material from outside lattice gauge theory into the Hamiltonian setting, each through an explicit typed map (Section~\ref{sec:side}): the incompressible flow of the Jacobian-conjecture counterexample becomes a family of local-energy trial states, whose energies track the free two-gluon threshold; a Navier--Stokes velocity field becomes a reducible $SU(2)$ connection, whose curvature masses grow at rates that the companion Navier--Stokes paper derives under four named statements of the imported manuscript \cite{NSPaper}; and the period data of the $S^6$ construction enter as a magnetic background through a single positive number. The zeta programme's Gram-sandwich lemma is used in the heat edition, and a Collatz identity serves as a regression fixture. In the author's view these connections are among the most valuable outcomes of the research; Section~\ref{sec:directions} gives my assessment of each.

\subsection*{How this record was made}

This record is an audit, carried out on 26--27 September 2026, of the workbench's published documents. The record's \texttt{README.md} says that the human participant proposed the research directions and possible connections and that Codex developed, checked and corrected the written constructions, and it attributes the public stewardship to KokunoYumeto, the owner of the workbench; the repository's README describes the workbench as open and AI-assisted, including work with ChatGPT~5.6 Sol and GPT-6 Astra. The workbench credits the exponential form of the vacuum and its character and Casimir calculus to Schütte, Zheng and Hamer \cite{SchutteZhengHamer}, the Fourier-algebra estimates to Eymard \cite{Eymard}, semigroup generation to Lumer and Phillips \cite{LumerPhillips}, and the creation method of the quantum line's first volume-uniform gap to Yarotsky \cite{Yarotsky}, with Baez's spin networks \cite{Baez} for the physical space. Its attribution file credits two inputs of side branches to Levent Alpöge: a polynomial map that the file identifies with the counterexample to the Jacobian conjecture announced by Alpöge on 20 July 2026, whose announcement, as the file records, credits Akhil for the question and Fable for the work; and the construction of a complex threefold circulated by Alpöge \cite{AlpogeS6}, which the file says was produced with Claude. The Navier--Stokes material rests on OpenAI's manuscript \cite{OpenAINS}, which the workbench cites and does not claim to have verified.

The audit was carried out by Claude (Anthropic), model \texttt{claude-opus-5-5} (\emph{claude-ab} below). It read the two proof manuscripts of the latest edition in full (F1--F42 and H1--H34), re-derived the analytic arguments, recomputed the constants and thresholds in exact arithmetic, recomputed the order-2 inputs from explicit $SU(2)$ tensors, and re-ran the workbench's order-5 producer and auditor, with byte-identical output. Independent Claude instances, each of which wrote none of the text, made a first-pass audit and referee passes; the referee pass of version~1 read for overstatement, understatement and missed generality, and several results stated here were found by it, each credited at its statement and verified by claude-ab before it was applied.

Every statement has one of four statuses:
\begin{itemize}
\item \emph{verified};
\item \emph{not verified here}, which implies nothing in either direction;
\item \emph{my assessment}, a view in the first person with its reasons, introduced by those words;
\item \emph{proved negative}, stated with its counterexample or derivation; these are collected in Section~\ref{sec:negative}.
\end{itemize}
The status lines under the results make the first two precise. \textsf{Proved here}, \textsf{Generalised here} (proved under weaker hypotheses or with a sharper constant than the record states), \textsf{Audited} (the workbench's proof read in full, every step checked, independent checks run) and \textsf{Reproduced} (the finite content recomputed by an independent program from the stated inputs) are forms of \emph{verified}. \textsf{Conditional} means proved given named inputs that were not checked here. \textsf{Audited in part} is verified in the parts named, and \textsf{Located} (found and described, not checked) is a form of \emph{not verified here}. The workbench itself separates written analytic arguments from finite exact-arithmetic certificates and from numerical diagnostics, and says so at each result.

The record was read at commit \texttt{fa79faf} of \texttt{KokunoYumeto/yang-mills-interacting-workbench} and in its 43 files (Section~\ref{sec:corpus}). Version~2 of this paper and the audit notes cited as \note{NN\_} were withdrawn from the branch \texttt{claude/claude-ab-grind-20260925} on 27 September 2026 because of their framing. They remain, unchanged, in the branch's history at commit \texttt{a90d9e5b}, in the folder \texttt{contrib/claude-ab/grind\_20260925/}, where the folders cited as \note{checks/\ldots} also are. The scripts on which this paper's statements rest are republished in \texttt{contrib/claude-ab/rewrite/checks/ym/} on the same branch.

This is version~3. It keeps the theorems, proofs and citations of version~2 and corrects its framing; it adds a section on the directions that remain open, with my assessment, and the typed maps of the side branches are stated with what each carries. Appendix~\ref{app:changes} lists every change.

\subsection*{Organization}

Section~\ref{sec:problem} states the problem, the setting and what the workbench claims, and Section~\ref{sec:corpus} maps the record. Section~\ref{sec:gap} proves the gap and Section~\ref{sec:heat} the heat and relaxation results. Section~\ref{sec:earlier} treats the earlier editions and Section~\ref{sec:side} the side branches and their typed maps. Section~\ref{sec:negative} collects the proved limits of specific routes, Section~\ref{sec:overlaps} the edges to the author's other workbenches, Section~\ref{sec:directions} the open directions with my assessment, and Section~\ref{sec:open} the questions and what was not checked. The appendices give the catalogue of statements, the verification, and the changes of version~3.
